\begin{exercisechunk}
\item \textbf{Backpropagation, affine scans, and recurrent attention.}
\label[exercise]{ex:l5-bptt}
Consider \(R_{t+1}=h_\theta(R_t,u_t)\) with observed inputs.
\begin{enumerate}[(a),beginpenalty=10000]
\item Express \(\partial R_T/\partial R_s\) as a product of
state-update Jacobians. For a scalar linear update \(R_{t+1}=aR_t\),
compare \(0<|a|<1\) and \(|a|>1\).
\item Verify associativity of the affine composition in
\cref{sec:rnn}. Compare a dense \(d\)-dimensional sequential recurrence
and a parallel scan in work, depth, and activation storage.
What improves if the transition matrices are diagonal?
\item Show that linearity in the state does not make learning the
transition parameter convex: use \(R_0=1\), \(R_2=a^2\),
and loss \((R_2-1)^2\).
\item Derive \cref{prop:attention-recurrence}. For \(\gamma_t=1\),
derive the extra denominator state needed for normalized linear
attention. State when the denominator is valid.
\end{enumerate}
\end{exercisechunk}
